Throughout, the gravitational field is the Painlev\'e-Gullstrand gauge field $\hb = 1 - \beta P$ of Fact \claimref{sheaf.gauge-field-cometric}, $\beta = \sqrt{2M/r}$, on flat spacetime (Definition \claimref{obj.synthesis.gtg-flat-spacetime-gauge-fields}), invertible on $r > 0$, with derived line element $dt^{2} - (dr + \beta\,dt)^{2} - r^{2}d\Omega^{2}$ read through Fact \claimref{synthesis.imported-gtg-dictionary}(i). The field does not depend on $t$, so the background translation $t \mapsto t + \tau$ is a symmetry of the gravitational field and $K = \hb^{-1}(e_{t})$ is a Killing vector in the sense of Fact \claimref{synthesis.imported-gtg-dictionary}(vi); $Q(K) = 1 - \beta^{2} = 1 - 2M/r$ vanishes exactly on $r = 2M$, where $K$ is normal, and the surface gravity there is $\kappa = 1/(4M)$, the value fixed by Theorem \claimref{singularity.euclidean-section-period} (imaginary Killing time of period $2\pi/\kappa = 8\pi M$ closes the Euclidean section) and recomputed in (ii) below as the non-affinity of the horizon generator. Write $s = \sqrt{r/2M}$, $\rho = s^{2} = r/2M$, so $\beta = 1/s$.

\emph{(i)} Radial null directions of the derived metric satisfy $(1 - \beta^{2})\dot t^{2} - 2\beta\,\dot t\dot r - \dot r^{2} = 0$, so $dr/dt = 1 - \beta$ (outgoing) or $dr/dt = -(1 + \beta)$ (ingoing); the labels $u = t - \int dr/(1 - \beta)$ and $v = t + \int dr/(1 + \beta)$ are constant along the outgoing and ingoing worldlines respectively. With $r = 2Ms^{2}$, $dr = 4Ms\,ds$ and $1 \mp \beta = (s \mp 1)/s$,
\[
\kappa\int\frac{dr}{1 \mp \beta} \;=\; \int\frac{s^{2}\,ds}{s \mp 1} \;=\; \int\Bigl(s \pm 1 + \frac{1}{s \mp 1}\Bigr)ds \;=\; \frac{s^{2}}{2} \pm s + \ln|s \mp 1| ,
\]
the partial fractions $s^{2}/(s - 1) = s + 1 + 1/(s - 1)$ and $s^{2}/(s + 1) = s - 1 + 1/(s + 1)$ being exact identities of rational functions (checked). Hence, for $s > 1$, $e^{-\kappa u} = e^{-\kappa t}(s - 1)e^{s^{2}/2 + s}$ and $e^{\kappa v} = e^{\kappa t}(s + 1)e^{s^{2}/2 - s}$, and the functions $U = -e^{-\kappa t}(s - 1)e^{s^{2}/2 + s}$, $V = e^{\kappa t}(s + 1)e^{s^{2}/2 - s}$ of the statement equal $-e^{-\kappa u}$ and $e^{\kappa v}$ there. Both are smooth on $r > 0$ ($s$ is smooth in $r$ there). $UV = -(s - 1)(s + 1)e^{s^{2}} = -(\rho - 1)e^{\rho}$, the exponents adding to $s^{2}$ and the prefactors to $s^{2} - 1$ (checked). $U$ vanishes exactly where $s = 1$, is negative for $s > 1$ and positive for $s < 1$, since $e^{(\cdot)}$ never vanishes; $V > 0$ for all $s \ge 0$. Differentiating a product $P(s)e^{E(s)}$ by $(P' + PE')e^{E}$: $\partial_{s}\bigl[(s - 1)e^{s^{2}/2 + s}\bigr] = \bigl(1 + (s - 1)(s + 1)\bigr)e^{s^{2}/2 + s} = s^{2}e^{s^{2}/2 + s}$ (checked), so $\partial_{s}U = -s^{2}e^{s^{2}/2 + s}e^{-\kappa t} \ne 0$ for $s > 0$. The symmetry $t \mapsto t + \tau$ multiplies $U$ by $e^{-\kappa\tau}$ and $V$ by $e^{\kappa\tau}$, because $t$ enters only through $e^{\mp\kappa t}$: $\partial_{t}U = -\kappa U$, $\partial_{t}V = \kappa V$, and $UV$ is invariant.

\emph{(ii)} In the chart $(v, r)$ the derived metric is $f\,dv^{2} - 2\,dv\,dr - r^{2}d\Omega^{2}$ with $f = 1 - 2M/r$ (substituting $dv = dt_{\mathrm{S}} + dr/f$ in $f\,dt_{\mathrm{S}}^{2} - dr^{2}/f$, $t_{\mathrm{S}}$ the Schwarzschild label $t - \int\beta\,dr/(1 - \beta^{2})$). The radial direction $\partial_{r}$ at fixed $v$ has $g(\partial_{r}, \partial_{r}) = 0$: the hypersurfaces $v = \mathrm{const}$, which are the level sets of $V = e^{\kappa v}$, are null, generated by the ingoing radial null worldlines. In that chart only $r$-derivatives of the metric are nonzero, and with $g^{vv} = 0$, $g^{vr} = -1$, $g^{rr} = -f$: $\Gamma^{r}{}_{rr} = \tfrac12 g^{rd}(2\partial_{r}g_{dr} - \partial_{d}g_{rr}) = 0$, since $g_{vr}$ is constant and $g_{rr} = 0$; so $r$ is an affine parameter along $v = \mathrm{const}$, and these worldlines are free-fall null worldlines by Fact \claimref{synthesis.imported-gtg-dictionary}(iii), crossing $r = 2M$ regularly. Likewise $\Gamma^{v}{}_{vv} = \tfrac12 g^{vr}(-\partial_{r}g_{vv}) = f'/2$, equal to $1/(4M)$ at $r = 2M$: the non-affinity of the horizon generator $\partial_{v}$, the surface gravity $\kappa$ (checked, with the tabulated value at two masses). Along $v = \mathrm{const}$, $u = v - 2r^{*}$ with $r^{*} = \int dr/f = r + 2M\ln|r/2M - 1|$, so $2\kappa r^{*} = \rho + \ln|\rho - 1|$ and $U = -e^{-\kappa u} = -e^{-\kappa v}(\rho - 1)e^{\rho}$, whose $\rho$-derivative is $-e^{-\kappa v}\rho\,e^{\rho}$ by the same product rule (checked): nonzero for $\rho > 0$ and equal to $-e^{1 - \kappa v}$ on the horizon. So $U$ is a smooth strictly monotone function of the affine parameter on each ingoing generator, with nonvanishing derivative at the horizon: a regular parameter. The level sets of $U$ are the hypersurfaces $u = \mathrm{const}$ generated by the outgoing radial null worldlines. Under the s-wave reduction premise (J3), Premise \claimref{hawking.joint-j3-s-wave-reduction}, each $(\omega, \ell, m, s)$ channel of a massless field on the exterior is a two-dimensional boundary mode, so the outgoing sector near the horizon is one chiral field per light ray, constant along the outgoing worldlines: a function of $U$ and the transverse labels, read off along any ingoing generator. This reduction is a hypothesis of the result, not derived here. Its carrier on the ingoing generator through $p \in \mathcal{H}^{+}$ is the continuum mode space $\CAR(L^{2}(\R, dU)\otimes K)$ of Definition \claimref{obj.hawking.continuum-mode-space}, the mode set attached to the sheet of $p$ generated by the ingoing null direction, in the labelling by null direction of Proposition \claimref{synthesis.r-loc-needs-sheet-labelled-modes}. By (i) the symmetry acts on it through $U \mapsto e^{-\kappa\tau}U$, and the weight that preserves the anticommutator $\{\psi(U), \psi^{*}(U')\} = \delta(U - U')$ is $\tfrac12$: $\beta_{\tau}(\psi(U)) = e^{-\kappa\tau/2}\psi(e^{-\kappa\tau}U)$, the boost of Theorem \claimref{synthesis.positive-null-energy-gives-r-premises} with $\xi = -\kappa U\partial_{U}$. On $U < 0$ (outside) the flow moves $U$ towards $0$, the future along the ingoing generator, so it is future-directed there; on $U > 0$ (inside) it moves $U$ away from $0$, the past. Both halves exist at the crossing for any hole whose late-time horizon lies in the Painlev\'e-Gullstrand field; the interior half ends at $r = 0$, where $\rho = 0$ and $U = e^{-\kappa v}$, so $U$ ranges over the half-line $(-\infty, e^{-\kappa v})$ on the generator, on which the translation group of (N$_{+}$) does not act (the generator of translations on $L^{2}$ of a half-line has no self-adjoint extension); the condition of (v) is therefore posed in the local form of (iv), whose limit object lives on the whole line.

\emph{(iii)} This is Theorem \claimref{synthesis.null-vacuum-singularity-under-reparametrisation-and-dilation}(i), where the argument (formerly given here) now stands: with $a = \xi'(U_{0})$, $b = \tfrac12\xi''(U_{0})$, $c = \tfrac16\xi'''(U_{0})$ the squared ratio of the weight-$\tfrac12$ kernels is $1 + \tfrac{ac - b^{2}}{a^{2}}x^{2} + O(x^{3}) = 1 + \tfrac16 S_{\xi}(U_{0})x^{2} + O(x^{3})$, so the kernels differ by $\tfrac{1}{12}S_{\xi}(U_{0})(U - U_{0}) + O((U - U_{0})^{2})$; weights $0$ and $1$ on the same kernel fail, and a weight-$h$ kernel has the continuous remainder $\tfrac{h}{6}S_{\xi}x^{2 - 2h}$. The check of this result still computes the squared ratio exactly at nine cases. The parameters $U$, $r$ and the affine parameter of (ii) are related by regular reparametrisations near the horizon, so the class of states with a bounded remainder is the same in all three.

\emph{(iv)} By (ii), $\beta_{\tau}(\psi(U)) = \varepsilon^{1/2}\psi(\varepsilon U)$ with $\varepsilon = e^{-\kappa\tau}$, the weight-$\tfrac12$ dilation $\delta_{\varepsilon}$ about $U = 0$, and $\tau \to \infty$ is $\varepsilon \to 0^{+}$; so (iv) is Theorem \claimref{synthesis.null-vacuum-singularity-under-reparametrisation-and-dilation}(ii), where the argument (formerly given here) now stands: after the substitution $U = \varepsilon x$ the singular part is exactly invariant and a bounded remainder contributes $\varepsilon\iint\bar f(x')g(x)R(\varepsilon x, \varepsilon x')\,dx\,dx' \to 0$, quasi-freeness (Fact \claimref{hawking.imported-quasi-free-kms-reduces-two}) carries the two-point limit to the whole carrier, the condition is sufficient and not necessary, and a remainder homogeneous of degree $-1$ survives unchanged. The ground state for the symmetry's own flow on the exterior half, the Boulware-type state, is invariant under $\beta_{\tau}$ and not equal to $\omega_{0}$ (its kernel is translation-invariant in $u = -\kappa^{-1}\ln(-U)$, not in $U$), so its remainder $W_{B} - W_{0}$ is $\beta_{\tau}$-invariant, that is homogeneous of degree $-1$ in $(U, U')$, and nonzero: it is not in the class.

\emph{(v)} At late times the ingoing generator through $p \in \mathcal{H}^{+}$ lies, on both sides of the horizon, in the region outside the collapsed matter, where the gravitational field is the Painlev\'e-Gullstrand field (the join to the matter's field is Birkhoff's theorem, literature); so (i) to (iv) apply on it. Under (N$_{+}^{\mathrm{coll}}$) the state of the outgoing modes on the sheet has kernel $W_{0} + R$ with $R$ bounded near $p$ in some regular parameter, hence in $U$ by (iii), and by (iv) its late-time limit under the symmetry's flow is $\omega_{0}$, the state of Theorem \claimref{synthesis.positive-null-energy-gives-r-premises}(i) for the translations in $U$. The exterior outgoing modes are those with support in $U < 0$, and the subalgebra $\CAR(L^{2}(\R_{-}, dU)\otimes K)$ is $\beta_{\tau}$-invariant. By Theorem \claimref{synthesis.positive-null-energy-gives-r-premises}(ii) the restriction of $\omega_{0}$ to the half-generator on which the boost at rate $\kappa$ is future-directed is the continuum I-state of Theorem \claimref{hawking.i-state-continuum} at $\beta = 2\pi/\kappa$ for that boost, jointly KMS, even and unique; here that half is $U < 0$ by (ii), and the flow $\xi = -\kappa U\partial_{U}$ on it is exactly the stored past-half clause of (ii), the boost $\beta_{-t}$ future-directed on $\CAR(L^{2}(\R_{-}, du))$ with symbol $(1 + e^{2\pi w})^{-1}$ on $|u|^{-1/2 + I_{4}w}$. That clause is a statement of CAR algebra on $L^{2}(\R, du)$ with the translations and the weight-$\tfrac12$ dilation, and nothing in its proof refers to which null line carries $u$, so it applies verbatim with $u = U$ on the ingoing generator, although the stored theorem names only the horizons (a) to (c). With $\kappa = 1/(4M)$, $\beta = 8\pi M$, and in Killing frequency $\varpi$ the occupation is $(1 + e^{2\pi\varpi/\kappa})^{-1}$, the symbol of Theorem \claimref{synthesis.positive-null-energy-gives-r-premises}(iii) (checked: the Mellin symbol at $w = \varpi/\kappa$ satisfies the two-point KMS identity at $\beta\kappa = 2\pi$ and at no other tested multiple). This is the late-time state of the near-horizon exterior outgoing modes of the collapse hole, thermal at $\kappa/2\pi = 1/(8\pi M)$ on the continuum carrier. That it is Hawking's flux at infinity uses Premise (I), Premise \claimref{hawking.premise-i-interface-identification}, which carries the occupancy of each near-horizon channel to the asymptotic channel by the exterior propagation, with the stored greybody inheritance; both are cited as hypotheses and not redone. The interior half $U > 0$ carries the restriction of $\omega_{0}$ for the reversed flow, the I-state at the same $\beta$ by the same theorem's past-half clause, and $\omega_{0}$ is pure on the whole generator (its symbol is a projection), so the exterior thermal state is the marginal of a pure state entangled across $\mathcal{H}^{+}$. The past half of $\mathcal{H}^{+}$, absent for a collapse hole, is never used: the half-generator cut is on the ingoing generator, at the horizon.

\emph{(vi)} The stored case (a) of Theorem \claimref{synthesis.positive-null-energy-gives-r-premises} poses (N$_{+}$) on the whole generators of the future horizon of the extended solution, and by the symmetry of the extended solution equally on those of the past horizon: a ground state on both sheets at the bifurcation surface, the Hartle-Hawking-type condition. (N$_{+}^{\mathrm{coll}}$) poses it on the ingoing sheet only, in the scaling limit at $\mathcal{H}^{+}$; the ingoing sector (functions of $V$) is not constrained by it and is fixed by the data of the collapse's past, which for a star collapsing from a region with vanishing gravitational field is the ground state in $v$ (no incoming flux): the Unruh-type condition. The Boulware-type state is excluded by (iv). The condition is sufficient, in the sense of (iii) and (iv): every state in the short-distance class has the ground state as its scaling limit; it is not necessary, since (iv) shows the same limit for every finite-rank change of the symbol. That a state prepared from regular data before the collapse is in the class at $\mathcal{H}^{+}$ is the propagation of short-distance regularity (Fulling, Sweeny and Wald 1978; the flux from it, Fredenhagen and Haag 1990), cited from memory and not imported: it is the justification the open problem asks for, argued here and not derived in the store; the companion result filed alongside shows it exactly for the in-vacuum of a null-shell collapse, whose ray-tracing map has a finite Schwarzian at the horizon ray.
